% Independent replacement fragments. Select sections per MERGE_GUIDE; not a standalone paper.

% Replace the prior-work transition in the Introduction, not the experimental results.
Prior designs already keep dirty data in the LLC while exposing writeback
candidates or dirty-state information separately~\cite{vwq,lastwrite,dbi}.
The question here is narrower: can a resident tag be reassigned while its
victim retains the same physical data slot until lower-level acceptance?
EP-L2 implements this ownership transition within a fixed payload and bank
budget. Its value must be established by separating copy removal, access to
otherwise unavailable staging capacity, and the service required to drain
that capacity.

% Replace the old Related Work section; do not append a second section.
% This is a literature module, not a new empirical claim or run authorization.
\section{Related Work}
\textbf{Writeback organization and scheduling.}
The Virtual Write Queue uses an LLC region to expose additional dirty-line
candidates to memory scheduling while retaining a physical write queue~\cite{vwq}.
Decoupled last-write prediction tracks predicted writebacks with address-only
entries while their data remain in the LLC~\cite{lastwrite}.
The Dirty-Block Index moves dirty-state tracking out of tag entries and groups
it by DRAM row~\cite{dbi}.
Thus, keeping dirty data in a cache or separating their metadata is not itself
new. EP-L2 studies a different transition: reusing a resident tag while the
old dirty payload remains in its original, independently owned slot, with
local data ownership and writeback ordering released at different events.

\textbf{Flexible data and transient-state storage.}
V-Way decouples tag and data placement for flexible associativity~\cite{vway};
Tag-Split Cache separates request and data-storage granularity~\cite{tsc}.
Unified GPU local memory reallocates storage among registers, cache, and
scratchpad~\cite{unified}.
In-cache MSHR organizations have a long history~\cite{farkas}; MiCache shares
FPGA storage between cache-line data and MSHR subentries~\cite{micache}.
These precedents motivate distinguishing the resource being shared, rather
than treating storage unification as the contribution. EP-L2 keeps the
resident-tag budget fixed and studies resident/WB payload ownership, not
extra lookup capacity or a shared data/request-metadata record.

\textbf{Miss handling and service constraints.}
DL-MSHR dynamically groups miss-tracking slots for primary and secondary
misses; its slots can also store write-miss data~\cite{dlmshr}.
Our primary comparison holds miss-tracking budgets fixed. Consequently,
reduced copy work and flexible staging must be evaluated separately from
additional request-tracking capacity and downstream service.

% Place near the downstream-headroom experiment rationale.
Request-buffer capacity affects scheduling visibility and implementation
complexity, not only queue depth~\cite{sms}. Moreover, reservation-failure
reductions need not be proportional to application speedup~\cite{dlmshr}.
We therefore distinguish buffering, service rate, and latency interventions,
compare B16 and U-WB under each identical condition, and report actual work
and elapsed cycles separately. These literature observations motivate the
controls; they do not establish the cause of our measured performance.

% OPTIONAL: retain only when victim retention / L1 request policies are discussed.
Victim caching retains displaced lines for subsequent reuse~\cite{victim},
whereas MRPB reorders and selectively bypasses requests before GPU L1
access~\cite{mrpb}. Neither is the primary WB-only comparison evaluated here;
any combination requires its own implementation and evaluation.
